\chapter{98}

\section{Mittwoch, 12.
September}



Manuela und Ava zogen noch vor Sonnenaufgang los. Vollgepackt mit all
dem Notwendigen.

«Wir stinken bald wie offene Müllsäcke, David und ich, und sind knapp am
Verhungern,» hatte Mazi am Vorabend hinunter gemeldet. Er hatte sich
gekonnt verzweifelt gegeben, aber darauf übermütig ins Handy gelacht.

Trotz Mazis humorvoller Einlage hatte sich Joh schäbig gefühlt.

«Wie kann es sein, B.B., dass ich mich so krass in Louis’ Drama
verliere? Manchmal fürchte ich, die beiden oben zu vergessen. Shit.»

Ben erhob sich vom Sofa. Mit Notizbuch und Stift näherte er sich Joh.
Der sass bereits am grossen Tisch und rückte eine Kaffeetasse auf seine
gegenüberliegende Seite. Ben setzte sich. Er fuhr sich mit den Fingern
durchs Haar und seufzte.

«Hm, ich such gerade nach Trost. Ich denk mir, grosse Katastrophen
generieren grosse Folgen.»

Dann richtete er sich gerade auf.

«Aber eigentlich ist es ein gutes Zeichen. Du vertraust darauf, dass
sie’s oben richten. Und das tun sie doch ganz gut, hm?»

«Auch wieder wahr, B.B.,» und Johs Stimme gewann an Bestimmtheit, «dann
also zu unserem Vorgehen. Sprechen wir unseren Plan nochmal durch, bevor
Louis da ist, okay?»



Unregelmässige Schrittgeräusche von oben näherten sich der Küche. Louis
kam humpelnd auf die beiden zu. Er wirkte geknickter als am Vorabend. Er
setzte sich neben Joh an den Tisch. Seine Niedergeschlagenheit war
unübersehbar.

Mit seiner eigenen Entlarvung hatte Louis zwar kurzfristig zur ersehnten
Beruhigung gefunden, aber gleichzeitig auch etwas verloren.

«Ich bin nicht Ciro. Es tut mir leid.»

Das unermüdliche Inszenieren einer Scheinwahrheit hatte sich mit dem
Aussprechen der beiden Sätze von einem Moment zum Nächsten erübrigt.
Damit war die Bombe geplatzt, und Louis’ weitere Worte hatten sein
Bekenntnis nur noch untermalt. Zusammen mit Ciro hatten ihn Motivation
und Kraft, als ein über die letzten Wochen eingeschworenes Helferduo,
sang- und klanglos verlassen.

Die Luft war raus.

Nun sass er vor ihnen. In diesem armseligen Zustand. Als Schuldiger, als
Feigling, gar Verräter, und unverfrorener Lügner. Da nützten ihm auch
Johs verständnisvolle Blicke und Bens Gelassenheit nichts. Die Last
hatte sich verlagert. Es ging nicht mehr darum, zu schauspielern, jetzt
ging es darum, sich zu stellen. Nicht nur den Menschen, die er belogen
hatte, auch sich selbst.
Wie ungefährlich die beiden dreinschauten. Es wäre passender und weit
einfacher, sie würden ihn anschreien. Ihre Milde war kaum zu ertragen.

«Packen wir’s an,» Joh rückte sich zurecht, «und noch mal, Louis, die
Vorkommnisse, die uns zwei betreffen, klären wir später, okay?»

Louis presste ein gequältes \emph{«Ja»} aus sich heraus, dem er ein
\emph{«Danke, Joh»} folgen liess. Er rutschte auf dem Stuhl hin und her,
wusste nicht, wohin mit seinen Armen und sah hilflos aus.

«Ich habe mich heute Morgen mit Ben beraten.»

Louis’ Augen begannen sofort wie Pingpong-Bälle zwischen Joh und Ben hin
und her zu hüpfen. Mit der Spannung wuchs gleichzeitig der Druck auf
seinen Magen. 

«Aber erstmal, Louis, konntest du dir bereits Gedanken machen?»

Joh legte den Kugelschreiber zur Seite. Über den Tisch hinweg berührte er 
Louis` Arm. Louis schaute überrascht hoch. Er wirkte, als hätte er die Frage nicht 
gehört oder nicht verstanden.
Joh neigte den Kopf zur Seite und suchte Louis’ Blick.

«Hej, Louis, du bist gerade weit unten. Das ist verständlich. Wir alle
können das nachvollziehen. Niemand von uns beneidet dich um deine Lage.
Wir helfen dir, okay?»

Johs Worte zeigten Wirkung. Sie lösten in Louis unvermittelt ein Gefühl
von Entspannung aus. Er richtete sich auf und wirkte augenblicklich
präsent.

«Ich bin euch sehr dankbar und echt erleichtert, Joh, dass du mich nicht
einfach rauswirfst. Es wär dir nicht zu verübeln.»

Louis legte unbeholfen seine Hand auf Johs, die nach wie vor auf seinem
Arm lag.
Dann wandte er sich an Ben.

«Und ehm, ob ich deine Unterstützung annehmen kann, Ben, ehm, ich weiss
es immer noch nicht. Wie soll ich dich bezahlen?», er löste sich von
Johs Hand und hob gestikulierend die Arme, «ich weiss nicht wie, wenn
ich an Anwaltskosten denke, dann, dann…»

Ben räusperte sich.

«Das weiss ich, Louis. Joh hat mich auch nach meinen Bedingungen
gefragt, solltest du einen Anwalt brauchen. Ich habe mir die Sache gut
überlegt und mache dir ein Angebot.»

«Ja,» Louis zögerte, «ja?»

«Was du auch immer an Unterstützung von mir brauchst, die nächsten Tage
voraussichtlich als Berater im Hintergrund, später, sofern notwendig, in
meiner Funktion als Anwalt, ich bin an deiner Seite. Ich tu das für
dich, und speziell, Louis, das sollst du wissen, für meinen Freund. Auch
Joh ist es ein Anliegen, dass sich die Sache ordentlich erledigt, im
Speziellen natürlich das Drumherum um deine Anstellung hier oben.»

Louis zuckte zusammen. Wieder kam das Gefühl der Beschämung hoch. Was
hatte er Joh nur angetan. Joh schien Louis’ Reaktion beobachtet zu
haben. Er nickte Louis aufmunternd zu.
Ben fuhr fort.

«Unter zwei Bedingungen bin ich solange für dich da, bis das Ganze
überstanden ist, und die müssen erfüllt sein. Erstens, Louis, du
\emph{musst} ehrlich sein zu mir. Ich muss jedes noch so kleinste Detail
des Unfalls kennen. Ich bin kein Richter, ich bin ein Verteidiger. Du
kannst mir die vollkommene Wahrheit zumuten und sie bleibt bei mir. Du
wirst unter meinem Schutz stehen. Keiner, nicht mal Joh, wird sie
erfahren,» Bens Blick war nun scharf auf Louis gerichtet, «und zweitens,
über jede deiner Handlungen, über alles, was du ab jetzt in dieser Sache
tust, musst du mich informieren, über ausnahmslos alles,» Bens Blick
wurde noch eindringlicher, «du hast bis morgen Zeit, eine Entscheidung
zu treffen. Ein Ja muss zu hundert Prozent aufrichtig sein, verstehst du
mich?»

Louis nickte überschnell, sah dabei konsterniert aus und kratzte sich an
der Schläfe. Es kam einem Wunder gleich. Ben würde ihn kostenlos
unterstützen? Er konnte es nicht fassen.
Joh lehnte sich entspannt zurück. Sein Ausdruck verriet, dass er Bens Bedingungen bereits kannte.


«Gut,» Ben blätterte in seinem Notizbuch, «auch wenn ich dein okay noch
nicht habe, wir müssen sofort deine nächsten Schritte planen, Louis, die
Zeit drängt. Uns bleiben drei Tage, die wir effizient nutzen müssen, um
eine Verschärfung der Lage zu vermeiden. Dass deine Freundin erpresst wird, 
ist heftig und entsprechend ernst zu nehmen. Joh und ich haben zusammen eine provisorische
Ablaufskizze erstellt. Ich schlage dir vor, dass wir zwei uns heute
Nachmittag an einen ungestörten Ort zurückziehen und du mir das Unglück
ganz genau beschreibst. Dann schauen wir, was wir als Erstes tun. Bist
du damit einverstanden?»

«Danke, Ben, danke, natürlich. Und, und deinen Bedingungen kann ich
schon jetzt zustimmen…»

«Gut, Louis. Trotzdem, ich bitte dich, die Entscheidung bis morgen genau
zu überdenken, darüber zu schlafen und sie mir morgen definitiv
mitzuteilen, okay?»

Louis nickte. Diesmal bedächtiger.
